---
date: 2026-06-17
tags:
  - marxism
  - essay
aliases:
---

up:: index

\section{Introdução}
Estamos no nível de abstração que Marx analisa no Livro I d'O Capital. Falamos da esfera de produção de valor, abstraindo de fatores como circulação e efetivação destes valores produzidos, i.e. assumindo que estes decorram como se não prejudicassem o movimento ``puro'' da produção. 

%%
\section{Aumento da eficiência, \textit{ceteris paribus}}
Um capital que possua maior eficiência em seu processo de produção consegue despender \textit{menos} meios de produção para um mesmo nível de produto. Portanto, comparado com as condições médias, ele consegue adiantar 
$$
(1-\sigma)c + v
$$

onde $\sigma>0$ indica o percentual a menos que ele consegue ``poupar'' de meios de produção em seu adiantamento de capital. Sua produção com este capital terá valor
$$
(1-\sigma)c + (1+m^{\prime})v
$$

%%

\section{Aumento de produtividade do trabalho, \textit{ceteris paribus}}
\begin{quote}
"Produtividade sempre diz respeito ao trabalho. O nível de produtividade nos informa uma determinada relação entre o trabalho vivo executado e o produto. (...)

O ganho de produtividade refere-se à diminuição do trabalho vivo exigido para produzir-se uma mercadoria (ou qualquer quantidade determinada de mercadorias). Como estão abstraídas [a princípio] variações de eficiência, o conjunto de matérias-primas e matérias auxiliares consumidos continua sendo o mesmo para qualquer nível dado de produção." \parencite[p. 233]{SaBarreto2021a}
\end{quote}

O aumento da produtividade do trabalho — chamado por Marx de ``forças produtivas''\footnote{Em alemão, \textit{Produktivkraft}; no plural, \textit{Produktivkräfte}.} — permite com que o dispêndio de força de trabalho para a produção de uma mercadoria seja reduzido. Ou seja, pode-se pensar que ``uma unidade'' de força de trabalho é capaz de mobilizar ``mais unidades'' de meios de produção. Dito de outra forma: a composição orgânica do capital $\frac{c}{v}$ aumenta.\footnote{A rigor, a composição \textit{orgânica} é uma unidade dialética da composição \textit{técnica} do capital — sua composição em meios de produção e de força de trabalho — e da composição \textit{de valor} do capital — esta sim igual a $\frac{c}{v}$. A menção de ``composição orgânica'' ao longo do texto — e, de fato, ao longo d'O Capital — diz respeito à composição de valor \textit{relacionada a composição técnica}, pois a análise isolada destas composições é perigosa: uma mesma composição técnica pode representar composições de valor diferentes ao longo do tempo, e vice-versa.}

Um exemplo numérico será de auxílio. Suponhamos que a produção a se analisar se dá em algum período de tempo $\Delta t$ dado, a uma taxa de exploração $m^{\prime}=100\%$; além disso, tenha-se que o capital médio adianta $c=v=50$ — logo, $c+v=100$ — em sua produção. Note-se que sua composição orgânica, portanto, é $\frac{c}{v}=100\%$.

Suponhamos agora que haja algum capital individual $K$ adiante também $100$, mas que possua um grau de produtividade $p=2$, ou seja, ele consegue mobilizar, por exemplo, $c=50$ adiantando somente $\frac{v}{p}=25$. Portanto, sua composição orgânica é $\frac{c}{\frac{v}{p}}=200\%$. 

Além disso, como ele consegue mobilizar $c=50$ meramente com $v=25$, ele agora tem um certo capital liberado $100-75=25$, a despender como lhe aprouver. Ele deseja, é claro, reinvesti-lo neste próprio processo, posto que já anotou em seu balancete que $100$ estaria investido nesta empreitada. Precisa, portanto, reinvestir este $25$ restante neste mesmo processo de produção, tomando em conta esta nova composição de capital. Ele desejará, portanto, investir uma certa quantidade $\tilde{c}$ em capital constante e $\tilde{v}$ em capital variável, de tal forma que $\tilde{c}+\tilde{v}=100$ e que $\tilde{c} = 2 \tilde{v}$ (diferente da média, para a qual $c=1v$). Fazendo umas contas, este capitalista chega à conclusão que consegue investir seu capital adiantado $100$ em $\tilde{c}=\frac{200}{3}$ e $\tilde{v}=\frac{100}{3}$.

Ou seja, o valor que ele produzirá será $\frac{200}{3} + \frac{100}{3} + \frac{100}{3}=100+\frac{100}{3}$, à taxa de exploração média $100\%$, enquanto o capital médio produzirá $50+50+50=150$. Claramente este capital produz um valor bem menor que a média, mas venderá seus produtos \textit{ao valor médio}; conseguirá obter, portanto, um mais-valor adicional, referente a esta diferença de valores efetivados, igual a $150-100 - \frac{100}{3}=50-\frac{100}{3}=\frac{50}{3}$. Ou seja, ele produziu um mais-valor menor que a média — $\frac{100}{3}<50$ —, porém compensou-o pela apropriação de $\frac{50}{3}$; no final das contas, este capital se apropriou de $\frac{100+50}{3} = 50$, i.e. do mesmo mais-valor que a média efetiva no mercado — o faz, porém, adiantando menos capital variável que a média!

Agora, tratemos formalmente este problema. Tenhamos que, em um dado período de tempo $\Delta t$, o capital médio despenda $c$ em capital constante e $v$ em capital variável, portanto adiantando um capital total $c+v$ e a uma composição orgânica igual a $\frac{c}{v}$, produzindo $N$ mercadorias a uma dada taxa de mais-valor $m^{\prime}$. Tenhamos também um capital específico $K$ que possui uma produtividade do trabalho maior que a média, de forma que, dado um capital constante $c$ adiantado, $K$ consiga despender uma quantidade $\frac{v}{p}<v$ de capital variável — onde $p>1$ mede o ``grau de produtividade'' deste capital com relação à média — para mobilizá-lo.

Note-se, portanto, que este capital não possui $\frac{c}{v}$ de composição orgânica, e sim $\frac{c}{\left( \frac{v}{p} \right)} = \frac{pc}{v}$, o que pode ser visto de duas formas: um capital variável $\frac{v}{p}<v$ é suficiente para mobilizar um capital constante $c$; e um capital variável $v$ consegue mobilizar capital variável $pc>c$.\footnote{A demonstração disso é por ``regra de três'': um capital variável $\frac{v}{p}$ está para $c$ assim como $v$ está para $pc$.}

Destaque-se que $K$ terá agora um $\Delta v$ ``a mais'' para reinvestir na produção\footnote{Neste caso, este capital a mais é igual a $v-\frac{v}{p} = \frac{p-1}{p}v < v$. Uma forma de entender o que essa fórmula significa é notar que ela diz que $v$ divide-se agora, para este capital mais produtivo, em uma parcela ``necessária'' $\frac{1}{p}v$ e uma parcela de valor ``liberado'' $\frac{p-1}{p}v$. }, tanto em capital constante como em capital variável. Evidentemente ele desejará reinvesti-lo, pois capital parado é uma contradição em termos — é um ``não-capital''; tem de investi-lo, porém, de acordo com sua (nova) composição orgânica. 

Para encontrar como ele investirá todo seu capital adiantado $c+v$ — já determinado \textit{ex ante} —, temos de notar que ele tem de satisfazer:
\begin{enumerate}
\item Ele deseja investir o mesmo capital adiantado $c+v$ em uma nova quantidade $\tilde{c}$ e $\tilde{v}$.\footnote{Note-se que $\tilde{c}>c$ e $\tilde{v}<v$. Contudo, destaque-se que $\frac{v}{p}< \tilde{v} < v$, porque $\frac{v}{p}$ é o capital variável necessário para mobilizar $c$, e $\tilde{v}$ é o necessário para mobilizar $c$ \textit{e também} uma parcela $\Delta c$ (tal que $\tilde{c} \equiv c + \Delta c$). Pode-se provar que $\tilde{v}$ tem de ser estritamente menor que $v$, porque $c+v=\tilde{c}+\tilde{v}$, o que faz com que $\tilde{v}= (c-\tilde{c})+v<v$, pois $\tilde{c}>c$ e, portanto, $(c-\tilde{c})<0$.  } 
\item Ele deseja reinvestir este capital de acordo com sua composição orgânica, que é maior que a composição média de acordo com seu grau de produtividade $p>1$. 
\end{enumerate}

Estas condições dão o sistema de equações:
$$
\begin{cases}
c+v&=\tilde{c}+\tilde{v} \\
\frac{\tilde{c}}{\tilde{v}} &= p \frac{c}{v}
\end{cases}
$$

Como $p$, $c$ e $v$ são dados, temos duas equações e duas variáveis $\tilde{c}$ e $\tilde{v}$ e, portanto, podemos obter a fórmula destas variáveis em função das informações dadas. Da segunda equação obtemos
$$
\tilde{c} = p \frac{c}{v} \tilde{v}
$$

Ressubstituindo na primeira equação, temos
\begin{align*}
\tilde{c}+\tilde{v} = \left( 1+p \frac{c}{v} \right) \tilde{v} = c+v
\end{align*}

Isolando $\tilde{v}$, e ressubstituindo para $\tilde{c}$, temos as soluções
$$
\begin{cases}
\tilde{v} = \frac{c+v}{1+p \frac{c}{v}} &\equiv \frac{1}{1+p \frac{c}{v}} C \\
\tilde{c} = p \frac{c}{v} \frac{c+v}{1+p \frac{c}{v}} &\equiv \frac{p \frac{c}{v}}{1+p \frac{c}{v}} C
\end{cases}
$$

onde $C$ é o capital adiantado $c+v$. Como \textit{sanity check}, pode-se checar que $\tilde{c}+\tilde{v} = C \equiv c+v$. Ou seja, este capital está empregando o mesmo capital adiantado $C$ que o capital médio emprega na produção de mercadorias em um período de tempo $\Delta t$.

Este capital $K$ somente adotará este novo ``plano de produção'' mais produtivo caso ele lhe dê uma vantagem contra as condições médias de mercado. Como Marx o descreve no capítulo de mais-valor relativo do livro I, $K$ obtém uma vantagem comparado à média quando o valor unitário — valor produzido dividido por cada mercadoria produzida — que consegue produzir é menor ou igual ao valor unitário médio\footnote{A rigor, ``valor médio'' é uma redundância, pois valor é naturalmente uma categoria \textit{social} (pode-se ver isso em que sua grandeza é tempo de trabalho \textit{socialmente} necessário). Faz sentido neste contexto — e também nas discussões de Marx no livro I — por se tratar dos valores \textit{produzidos} por um capital individual e pelo capital médio.}. 

Temos que o capital médio produz valor
$$
c + (1+m^{\prime})v
$$

e o capital $K$ produz
$$
\tilde{c}+(1+m^{\prime})\tilde{v}
$$

\textit{Ex hypothesi}, o capital médio produz $N$ mercadorias em um período de tempo $\Delta t$ dado; digamos que este capital mais produtivo $K$ produza $N+\Delta N$ mercadorias, onde $\Delta N \geq 0$. Então, a relação entre seus valores unitários pode ser obtida da seguinte forma: 
$$
\frac{\tilde{c} + (1+m^{\prime})\tilde{v}}{N+\Delta N} < \frac{c+(1+m^{\prime})v}{N+\Delta N} \leq \frac{c+(1+m^{\prime})v}{N}
$$

De fato, um capital mais produtivo \textit{garantidamente} produz mercadorias com valor unitário menor que o da média, \textit{mesmo quando não produz mercadorias a mais}: como $c+v = \tilde{c}+\tilde{v}$, então, se ele produz a mesma quantidade $N$ de mercadorias, ele só tomaria vantagem se $m^{\prime}\tilde{v} \leq m^{\prime}v$\footnote{Pois os denominadores dos valores unitários seriam iguais.} — o que é satisfeito justamente por seu grau maior de produtividade, que faz com que $\tilde{v} < v$. 

Ou seja, com um aumento do grau de produtividade, ele \textit{sempre} consegue adquirir um mais-valor adicional $m^{\prime}(v-\tilde{v}) > 0$, pois esta é justamente a diferença entre o valor produzido pela média e o valor produzido por $K$. Portanto, abrindo $\tilde{v}$, e chamando seu mais-valor adicional de $\Delta m$, temos
$$
\Delta m = m^{\prime}(v-\tilde{v})
$$

Como temos que 
$$
\tilde{v} = \frac{c+v}{1+p \frac{c}{v}} = \frac{\left( 1+\frac{c}{v} \right)}{1+p \frac{c}{v}} v
$$

obtemos que
\begin{align*}
\Delta m &=\frac{\left( 1 + p \frac{c}{v} - 1 - \frac{c}{v} \right)}{1+p \frac{c}{v}} m^{\prime}v \\
&= \frac{(p-1) \frac{c}{v}}{1+ p \frac{c}{v}} m^{\prime}v
\end{align*}

(Como \textit{sanity check}, veja que $p=1$ faz com que $\Delta m=0$, como esperado.)

Ademais, seu mais-valor \textit{produzido} $\tilde{m}$ é 
\begin{align*}
\tilde{m} &= m^{\prime}\tilde{v}  \\
&= m^{\prime}v - \Delta m \\
&= \frac{\left( 1+\frac{c}{v} \right)}{1 + p \frac{c}{v}} m^{\prime} v  \\
&< m^{\prime}v
\end{align*}

Ou seja, enquanto o capital médio obtém mais-valor $m^{\prime}v$ — que ele próprio produz e de que também se apropria —, o capital mais produtivo \textit{produz} mais-valor $m^{\prime} \tilde{v}$ \textit{menor} que $m^{\prime}v$, mas também se \textit{apropria} de mais-valor adicional $\Delta m$, de tal forma que, no fim e ao cabo, este capital se apropriará do mesmo mais-valor $m^{\prime}v$ que a média. Ou seja, $K$ produz menos mais-valor por si próprio, mas \textit{efetiva} um mais-valor adicional de tal magnitude que precisou adiantar \textit{menos} força de trabalho (do que a média) para, mesmo assim, obter o \textit{mesmo} excedente de trabalho (i.e. mais-valor) como se tivesse produzido nas condições médias.\footnote{O fato de que seu mais-valor total é igual ao mais-valor que a média produz é consequência de que este capital \textit{adianta o mesmo capital que a média}, $c+v$. Já aqui podemos ver que adiantar o mesmo capital \textit{não} significa que suas apropriações de valor serão iguais.}

Pelas fórmulas de $\tilde{c}$ e $\tilde{v}$, fica clara a relação entre estas novas quantidades investidas e a composição orgânica maior deste capital. Neste nível de abstração, porém, estamos falando puramente da \textit{produção} de valor deste capital, em que a apropriação de valor deste capital vem \textit{do valor social} ao qual ele vende suas mercadorias. Porém, não é possível ver, neste nível, que o mais-valor adicional de que ele se apropria \textit{é produzido por capitais menos produtivos}, conclusão à qual podemos chegar somente no nível de abstração de taxas de lucro (que Marx aborda somente no livro 3 d'O Capital). 

Além disso, a taxa de mais-valor deste capital mais produtivo também se diferencia da taxa de mais-valor do capital médio, no sentido de que, assim como este capital \textit{produz} $\tilde{m}$ mas \textit{se apropria} de $m = \tilde{m} + \Delta m$, também sua taxa de mais-valor se diferencia entre aquela que advém de sua \textit{produção} e aquela que advém de sua \textit{apropriação}. A que diz respeito à sua produção é
$$
\frac{\tilde{m}}{\tilde{v}} = m^{\prime} \frac{\tilde{v}}{\tilde{v}} = m^{\prime}
$$

i.e. é, quase por definição, a mesma taxa de exploração que a da média. A taxa que diz respeito à sua \textit{apropriação}, porém, é diferente. Denotando-a por $m^{\dagger}$, temos que ela é
\begin{align*}
m^{\dagger} &= \frac{m^{\prime}v}{\tilde{v}}\\
&= \frac{1+p \frac{c}{v}}{1+\frac{c}{v}} m^{\prime} \\
&> m^{\prime}
\end{align*}

Ou seja, no que tange ao excedente que este capital \textit{de fato} se apropria, sua taxa de mais-valor \textit{é} maior que a da média; não é meramente uma aparência ``\textit{as-if}'' deste fenômeno, mas algo que é real: este capital adianta capital variável $\tilde{v}$, que reproduz seu próprio valor $\tilde{v}$ e um excedente/mais-valor $m^{\prime}\tilde{v}$, \textit{e também} apropria-se, após a venda de suas mercadorias pelo valor social, de um extra $\Delta m$. Não é mera aparência: ele, de fato, apodera-se de um valor excedente ao que adiantou previamente. Ademais, não só sua taxa de exploração aparece aos demais capitalistas como maior; o que eles enxergam, concretamente, são os preços menores que este capital é capaz de sustentar no mercado. Mesmo que estejamos aqui pressupondo (como Marx no livro I) que mercadorias são vendidas por seus valores, não só o capital mais produtivo produzirá uma quantidade maior de mercadorias, como o próprio valor ``bruto'' por ele produzido será menor; portanto, o valor unitário de cada mercadoria terá ``dois motivos'' para ser menor que o valor unitário da média: o ``valor bruto'' produzido é menor do que o que a média produz, e o valor unitário será menor ainda do que o valor unitário médio.

Em suma: o aumento de produtividade de um capital, \textit{vis-à-vis} produtividade média de seu setor, permite que ele 1) aumente seu emprego de capital constante, trabalho morto, e 2) diminua sua necessidade de capital variável, trabalho vivo, ao mesmo tempo em que 3) se apropria do mesmo excedente de que se apropriava antes de seu ganho de produtividade. Implícito está, de forma mistificada, que haverá valor ``disponível'' no mercado de que ele conseguirá se apropriar, i.e. que ele \textit{conseguirá efetivar} seu produto ao valor social (em nível mais concreto: ao preço de mercado), que haverá demanda para a oferta aumentada (em particular devido a este capital) de seu setor etc. Porém, ao passo que este capital é capaz de se apoderar de parcelas maiores de mercado — portanto, acumulando capital —, seu objetivo não é mais meramente conquistar a demanda de seus compradores, mas também, e talvez até mais, apoderar-se do capital de seus competidores,
\begin{quote}
``...concentração de capitais já constituídos, supressão [\textit{Aufhebung}] de sua independência individual, expropriação de capitalista por capitalista, conversão de muitos capitais menores em poucos capitais maiores.'' \parencite[p. 701]{Marx2017}
\end{quote}

O valor da força de trabalho diz respeito ao valor total das mercadorias demandadas por seus ofertantes — a classe trabalhadora assalariada — para sua reprodução dia após dia; não é porque seu consumo siga o ciclo $M-D-M'$, porém, que estes trabalhadores não consumam mercadorias produzidas por capitalistas — portanto, mercadorias cujos valores contêm em si mais-valor. A produtividade dos setores que produzem estas mercadorias, portanto, ao abaixarem os valores das mercadorias neles produzidas, rebaixam também o valor da força de trabalho dos trabalhadores que as consomem. Está claro, portanto, que o mecanismo que move o capitalista individual a aumentar sua produtividade do trabalho, em particular nestes setores de subsistência, é o mesmo mecanismo que induz o paulatino aumento da produtividade dos setores em que pertencem. Este ímpeto do capitalista particular, individualmente e contra a produção social à qual remete, é o mesmo que induz a queda do valor social, cuja diferença \textit{vis-à-vis} sua produção individual é o que lhe concede vantagem contra seus competidores, para começo de conversa; aquilo que lhe permite correr mais rápido que os demais capitalistas é o mesmo movimento que, quando perpetrado socialmente, faz com que sua liderança se esvaia pouco a pouco, induzindo com que ele aumente novamente sua capacidade produtiva — que tenha, portanto, não só um ímpeto crescente, como também \textit{acelerado}, crescentemente crescente. Dessa forma, é sempre efêmera sua sensação de estar ``ganhando do mercado'', sendo esta sensação, portanto, mera mistificação de que sua glória de Prometeu é tão maior quanto mais chegam águias, tanto em maior quantidade, quanto maior voracidade para devorar suas entranhas. 

%%
\section{Aumento da intensidade do trabalho, \textit{ceteris paribus}}
\begin{quote}
``No caso de aumento da intensidade, as novas tecnologias contribuem no sentido de concentrar uma quantidade maior de trabalho em um dado período de tempo. Lembrando que abstraímos de outros efeitos, trata-se da elevação do ritmo da produção sem que se alterem as razões entre o trabalho e produto (a produtividade) e entre matérias-primas [meios de produção] e produto (a eficiência). A maior intensidade atua, portanto, no sentido de comprimir uma jornada de trabalho mais extensa em uma jornada menor.'' \parencite[p. 234]{SaBarreto2021a}

``Ao lado da medida do tempo de trabalho como 'grandeza extensiva'[,] apresenta-se agora a medida de seu grau de condensação. A hora mais intensa da jornada de trabalho de 10 horas encerra tanto ou mais trabalho, isto é, força de trabalho despendida, que a hora mais porosa da jornada de trabalho de 12 horas. Seu produto [da 'hora mais intensa'] tem, por isso, tanto ou mais valor que o produto da $1 ^1/_{5}$ [$\frac{12}{10}$] hora mais porosa.'' \parencite[p. 482-3]{Marx2017}
\end{quote}

Um aumento puro de intensidade do trabalho ocorre quando há uma ``condensação'' do trabalho despendido. Isto se dá principalmente quando são estabelecidas leis trabalhistas que restringem jornadas de trabalho, em cujo caso o capital busca fazer com que seu capital adiantado renda-lhe o mesmo que rendia antes, nem que faça-o em menos tempo.
%%

---
\subsubsection{Referências}
